\documentclass[a4paper, 10pt]{article}

\usepackage[margin=1in]{geometry}
\usepackage{amsmath, amssymb}
\usepackage{mathpartir}
\usepackage{xcolor}
\usepackage{booktabs}
\usepackage{array}

\newcommand{\vof}[2]{(\texttt{value-of}\ #1\ #2)}
\newcommand{\ttt}[1]{\texttt{#1}}
\newcommand{\clos}[3]{\mathrm{clausura}(#1, #2, #3)}
\newcommand{\err}{\mathrm{error}}

\title{Tarea 2 \\ Unidad 2: Procedimientos y clausuras}
\author{Angel Fernando Bórquez Guerrero \\ Lenguajes de Programación 2026-2}
\date{}

\begin{document}
\maketitle

%% ------------------------------------------------------------------
\section*{1.1 Varios parámetros en el núcleo}

\begin{mathpar}
\inferrule*[right={[proc]}]
  { }
  {\vof{(\ttt{proc-exp}\ xs\ c)}{\rho} = \clos{xs}{c}{\rho}}
\end{mathpar}

\begin{mathpar}
\inferrule*[right={[call]}]
  {\vof{e_1}{\rho} = \clos{xs}{c}{\rho'} \\
   \vof{a_i}{\rho} = v_i \ \ (1 \leq i \leq n) \\
   |xs| = n}
  {\vof{(\ttt{call-exp}\ e_1\ es)}{\rho} =
   \vof{c}{[xs=\langle v_1, \dots, v_n\rangle]\rho'}}
\end{mathpar}

\begin{mathpar}
\inferrule*[right={[call-err]}]
  {\vof{e_1}{\rho} = \clos{xs}{c}{\rho'} \\
   \vof{a_i}{\rho} = v_i \ \ (1 \leq i \leq n) \\
   |xs| \neq n}
  {\vof{(\ttt{call-exp}\ e_1\ es)}{\rho} = \err}
\end{mathpar}

El error de [call-err] es
\ttt{apply-procedure: número de argumentos equivocado: esperaba $|xs|$ y me diste $n$}.

%% ------------------------------------------------------------------
\section*{1.2 Varios parámetros como azúcar}

\[
\begin{array}{lcll}
\ttt{proc ($x$, $y$, $ys$ ...) $c$} & \leadsto & \ttt{proc ($x$) proc ($y$, $ys$ ...) $c$} & \text{[proc*]} \\[4pt]
\ttt{($e_1$ $e_2$ $e_3$ $es$ ...)} & \leadsto & \ttt{(($e_1$ $e_2$) $e_3$ $es$ ...)} & \text{[call*]}
\end{array}
\]

Los argumentos se pasan de izquierda a derecha: el primer operando va al primer
\ttt{proc}.

%% ------------------------------------------------------------------
\section*{1.3 Comparación}

\begin{center}
\begin{tabular}{l>{\raggedright}p{6.2cm}>{\raggedright\arraybackslash}p{6.2cm}}
\toprule
 & \textbf{Núcleo} (\ttt{proc-varios-nucleo}) & \textbf{Azúcar} (\ttt{proc-varios-azucar}) \\
\midrule
P1 & \ttt{7} & \ttt{7} \\
P2 & \ttt{número de argumentos equivocado: esperaba 2 y me diste 1} & \ttt{7} \\
P3 & \ttt{número de argumentos equivocado: esperaba 2 y me diste 3} & \ttt{val->proc: Esperaba un procedimiento pero me diste 7} \\
P4 & \ttt{5} & \ttt{parse: no esperaba \textasciigrave)\textasciigrave{} después de \textasciigrave(\textasciigrave} \\
P5 & \ttt{número de argumentos equivocado: esperaba 1 y me diste 2} & \ttt{7} \\
\bottomrule
\end{tabular}
\end{center}

\paragraph{1. ¿Qué es \ttt{(resta 10)}?}
En el núcleo es un error: \ttt{resta} pide exactamente dos argumentos (P2). En el azúcar
es un procedimiento, \ttt{<proc (y)>}, la clausura del \ttt{proc} interior con $x=10$,
que espera el segundo argumento (P2 da \ttt{7}).

\paragraph{2. ¿Cuándo se detecta un operando de más?}
En el núcleo, en la llamada misma, antes de evaluar el cuerpo, y el mensaje dice cuántos
esperaba y cuántos recibió (P3). En el azúcar, \ttt{(resta 10 3 1)} se traduce a
\ttt{(((resta 10) 3) 1)}: se calcula $10-3=7$ y el error llega una llamada después,
al intentar llamar a \ttt{7}. El mensaje no habla de argumentos (P3).

\paragraph{3. ¿Qué escribe el REPL para \ttt{proc (x, y) x}?}
El núcleo escribe \ttt{<proc (x, y)>}, porque la clausura guarda la lista de
parámetros. El azúcar escribe \ttt{<proc (x)>}, porque \ttt{value-of} nunca ve un
\ttt{proc} de dos: recibe \ttt{proc (x) proc (y) x}, y su valor es la clausura del
exterior, que solo tiene a \ttt{x}.

\paragraph{Qué habilita y qué sacrifica cada uno.}
\begin{itemize}
  \item \textbf{El núcleo habilita} procedimientos sin parámetros (P4 vale \ttt{5}) y
    errores de aridad en el lugar exacto y con un mensaje claro (P3) Además, el valor
    dice la verdad sobre sus parámetros (pregunta 3).
  \item \textbf{El núcleo sacrifica} la aplicación parcial (P2) y la mezcla con
    procedimientos currificados: un \ttt{proc (x) proc (y)} ya no se puede llamar con dos
    operandos (P5).
  \item \textbf{El azúcar habilita} la aplicación parcial gratis (P2), y los
    procedimientos currificados y los de varios parámetros son intercambiables (P5 vale
    \ttt{7}). No hay que tocar el intérprete: las reglas de inferencia y las pruebas de
    \ttt{value-of} siguen igual.
  \item \textbf{El azúcar sacrifica} los procedimientos de cero parámetros, que no tienen
    traducción a la currificación (P4 ni siquiera se parsea). También sacrifica los
    errores de aridad, que llegan tarde y disfrazados (P3), y la información de los
    parámetros en el valor (pregunta 3).
\end{itemize}
Es un intercambio genuino: lo que en un enfoque es un error (P2, P5), en el otro es una
capacidad, y viceversa (P3, P4).

\paragraph{Lo que se modificó.}
\emph{Núcleo:} seis archivos. En \ttt{ast.rkt}, los campos \ttt{xs} y \ttt{es} con sus
contratos. En \ttt{parser.rkt}, las dos producciones y los no terminales
\ttt{parameters}, \ttt{nonempty-parameters} y \ttt{operands}. En \ttt{desugar.rkt}, la
reconstrucción con \ttt{(map desugar es)}. En \ttt{procs.rkt}, la clausura con
\ttt{xs}. En \ttt{interp.rkt}, las cláusulas de \ttt{proc-exp} y \ttt{call-exp}, y
\ttt{apply-procedure}, que comprueba la aridad y extiende con \ttt{extend-env*}. En
\ttt{vals.rkt}, \ttt{expval->string}. Es decir, todas las capas del intérprete.

\emph{Azúcar:} tres archivos. En \ttt{ast.rkt}, los \ttt{struct} \ttt{proc*-exp} y
\ttt{call*-exp}. En \ttt{parser.rkt}, las dos producciones, con
\ttt{nonempty-parameters} y \ttt{nonempty-operands}. En \ttt{desugar.rkt}, las
cláusulas de [proc*] y [call*]. Solo cambia el frente del intérprete: \ttt{value-of},
\ttt{procs.rkt} y \ttt{vals.rkt} quedan intactos.

%% ------------------------------------------------------------------
\section*{2.1 Vinculación dinámica: especificación}

\begin{mathpar}
\inferrule*[right={[proc]}]
  { }
  {\vof{(\ttt{proc-exp}\ x\ c)}{\rho} = \mathrm{proc}(x, c)}

\inferrule*[right={[call]}]
  {\vof{e_1}{\rho} = \mathrm{proc}(x, c) \\
   \vof{e_2}{\rho} = v_2}
  {\vof{(\ttt{call-exp}\ e_1\ e_2)}{\rho} = \vof{c}{[x=v_2]\rho}}
\end{mathpar}

Cambia la conclusión de [call]: el cuerpo se evalúa en $[x=v_2]\rho$, el entorno de la
llamada, y no en $[x=v_2]\rho'$. Como $\rho'$ ya no aparece en ninguna conclusión, el
valor de un \ttt{proc} no necesita guardar un entorno: es $\mathrm{proc}(x, c)$, y por
eso cambian también la conclusión de [proc] y la primera premisa de [call]. En el código,
la clausura perdió su campo $\rho$ (\ttt{procs.rkt}, \ttt{vals.rkt}), y
\ttt{apply-procedure} recibe el $\rho$ de la llamada (\ttt{interp.rkt}).

%% ------------------------------------------------------------------
\section*{2.3 Análisis}

\begin{center}
\begin{tabular}{l>{\raggedright}p{6.2cm}>{\raggedright\arraybackslash}p{6.2cm}}
\toprule
 & \textbf{PROC} (\ttt{proc-02-procedimientos}) & \textbf{Dinámica} (\ttt{proc-dyn}) \\
\midrule
D1 & \ttt{3} & \ttt{5} \\
D2 & \ttt{3} & \ttt{2} \\
D3 & \ttt{-100} & \ttt{0} \\
D4 & \ttt{7} & \ttt{apply-env: la variable n no está vinculada} \\
D5 & \ttt{apply-env: la variable y no está vinculada} & \ttt{5} \\
\bottomrule
\end{tabular}
\end{center}

\paragraph{1. Renombrar un parámetro (D1 y D2).}
La \ttt{a} del cuerpo de \ttt{p} es libre. En PROC se busca en el entorno de creación de
\ttt{p}, donde $a=3$, y nada de lo que pase después lo cambia: D1 y D2 valen \ttt{3}. Con
vinculación dinámica se busca en el entorno de la llamada \ttt{(p 0)}, que está dentro del
cuerpo de \ttt{f}. En D1 ese entorno tiene $x=2$ sobre $a=5$, y da \ttt{5}. En D2 el
parámetro de \ttt{f} se llama \ttt{a}, así que la \ttt{a} más cercana es $a=2$, y da
\ttt{2}. El nombre del parámetro de \ttt{f}, que debería ser asunto privado de \ttt{f},
se filtra al cuerpo de \ttt{p}. Quien escribe \ttt{f} tiene que conocer las variables
libres de todo procedimiento que llame, y quien escribe \ttt{p} tiene que conocer los
nombres vinculados en todos los lugares desde donde se le llame. Con alcance léxico basta
leer el texto donde está escrito \ttt{p}.

\paragraph{2. Un procedimiento que devuelve otro (D4).}
En PROC, \ttt{(menos 5)} es una clausura que guarda $n=5$, así que sigue funcionando
cuando \ttt{menos} ya terminó, y da \ttt{7}. Con vinculación dinámica,
\ttt{(menos 5)} devuelve $\mathrm{proc}(x, \ttt{-(x, n)})$ sin entorno. El vínculo
$n=5$ solo existió mientras corría el cuerpo de \ttt{menos}, y cuando se llama con
\ttt{12} el entorno de esa llamada no tiene \ttt{n}: es un error. La currificación
deja de funcionar, y con ella el azúcar de la parte 1.2.

\paragraph{Qué habilita y qué sacrifica cada diseño.}
\begin{itemize}
  \item \textbf{La vinculación dinámica habilita} que quien llama configure al
    procedimiento sin pasarle argumentos: en D5, \ttt{f} usa una \ttt{y} que se define
    después y vale \ttt{5}, como un parámetro implícito o una variable de ambiente. La
    representación también es más sencilla, porque un procedimiento no guarda entorno.
  \item \textbf{La vinculación dinámica sacrifica} la seguridad al renombrar (D1 y D2),
    los procedimientos que devuelven procedimientos (D4) y la posibilidad de saber qué
    hace un procedimiento leyendo su texto: en D3, \ttt{f} y \ttt{g} se escribieron en
    lugares distintos y se comportan igual, porque las dos ven $x=100$, y dan \ttt{0}.
  \item \textbf{El alcance léxico habilita} las clausuras (D4), que el significado no
    dependa de quién llama (D1, D2 y D3), y detectar errores sin correr el programa: en
    D5, \ttt{y} es libre en \ttt{f} y no hay ningún vínculo a la vista en el texto.
  \item \textbf{El alcance léxico sacrifica} la posibilidad de pasar valores desde la
    llamada de forma implícita (D5 es error) y obliga a guardar un entorno en cada
    procedimiento.
\end{itemize}

\paragraph{3. ¿Por qué el alcance léxico es la norma?}
Por la modularidad. Con vinculación dinámica, una función de biblioteca que use una
variable libre cambia de resultado según cómo nombre sus variables locales quien la llama,
como en D2, y nadie puede razonar sobre ella por separado. Además, las funciones de orden
superior (callbacks, \ttt{map} con una \ttt{lambda}, la currificación) necesitan
clausuras, y con vinculación dinámica se rompen como en D4. El alcance léxico también
permite que un compilador resuelva cada variable antes de correr y avise de errores, como
el \ttt{g++ -Wall} de la sesión 1. No es que lo dinámico no sirva: los lenguajes que lo
necesitan lo ofrecen de forma controlada y explícita (las variables especiales de Common
Lisp, \ttt{parameterize} en Racket, las variables de ambiente), y dejan el alcance léxico
como regla general.

\end{document}
